\section{Conclusion and Future Work}\label{sec:disc}

IoT deployments that keep inference on site run LLMs as pipelines over
heterogeneous, memory-constrained devices, where a stage failure loses the
request's KV state and existing recovery paths trade prefix-length
recomputation, sum-of-stages replication, or minutes of reconfiguration.
\sys{} recovers the lost KV from one or two cross-stage parity columns
without recomputing the prefix and selects the layer placement together with
the backup assignment so that the failed stage's weights have a device to
load on.
On a six-stage Jetson pipeline serving OPT-6.7B, single parity recovers in
1.2--1.4\,s, flat over the measured failure positions, while recomputation
grows to 16\,s.  It recovers about one decode step later than replication
while retaining $3.7\times$ less coordinator KV state ($2.7\times$ over
the stages the parity path recovers), and it changes failure-free throughput
by $+0.6$\% ($\pm11$\% across three rounds).  Double parity restores two
simultaneous stage failures from the same stored bytes at the cost of one
when their backups sit on distinct nodes.  In an offline placement analysis
with backups confined to pipeline devices, joint placement remains feasible at every
per-device memory cap at which the model fits, where a cost-first placement
cannot be backed up.  For edge deployments that cannot depend on a remote
service, this turns a protected-stage failure whose recovery contract holds
from a request loss or a minutes-long reconfiguration into a delay of two to
four decode steps.

Future work includes a full Reconfigure sweep at 7B and end-to-end
measurements that include heartbeat detection, process-kill injection, and
the lazy-backup middle point, which would complete the latency comparison,
and adding backup spread to the placement objective, which would make the
two-failure cost independent of where backups land.  Other model families, fleets with
different SKUs, and a head/middle/tail victim sweep would test the
generality of the flat-in-$P$ result.  Protecting the coordinator and the
head stage are the next steps for the protection domain.  A general
$k$-parity construction with Cauchy coefficients pays off only when
$\sum_s\ell_s/\ell_{\max}$ is large, and the contract can be relaxed for
tail victims and prefill failures while the placement objective absorbs
recovery latency and per-hop costs.
